\subsection{DML -- Data Manipulation Language}
Die DML umfasst alle Befehle, die Daten \emph{innerhalb} bestehender Tabellen verändern: einfügen (\icode{INSERT}), ändern (\icode{UPDATE}), ersetzen (\icode{REPLACE}) und löschen (\icode{DELETE}).

\subsubsection{INSERT}
\begin{syntaxbox}[SQL]
-- Mit Spaltenliste (empfohlen), mehrere Datensaetze auf einmal
INSERT INTO Leistung (LNr, Bezeichnung, Kosten, AID) VALUES
    (800, 'FamilienRecht',  67.00, NULL),
    (801, 'ArbeitsRecht',  700.00, NULL);

-- Ohne Spaltenliste: nur wenn ALLE Spalten in der Reihenfolge
INSERT INTO Leistung VALUES (804, 'Sonstiges', 99.00, NULL);
\end{syntaxbox}

\begin{infobox}[Hinweis zur Spaltenliste]
Die runden Klammern nach dem Tabellennamen (die Spaltenliste) kann man weglassen, wenn man alle Spalten in der korrekten Reihenfolge angibt. Die Klammern bei \icode{VALUES (...)} bleiben immer.
\end{infobox}

\subsubsection{UPDATE, REPLACE und DELETE}
\begin{syntaxbox}[SQL]
-- Eine oder mehrere Spalten aendern
UPDATE Leistung
SET AID = 2, Kosten = 99.00
WHERE LNr = 17;

-- Bestehenden Datensatz komplett ersetzen
REPLACE kurs_db.Abteilung VALUES (4, 'Rechnungswesen');

-- Datensaetze mit Bedingung loeschen
DELETE FROM Leistung
WHERE LNr = 800;
\end{syntaxbox}

\begin{warnbox}[Achtung]
\icode{UPDATE} und \icode{DELETE} ohne \icode{WHERE} verändern bzw. löschen \emph{alle} Zeilen der Tabelle. Vor dem Ausführen immer die \icode{WHERE}-Bedingung prüfen.
\end{warnbox}
